\documentclass[12pt]{article}
\usepackage[utf8]{inputenc}
\usepackage[T1]{fontenc}
\usepackage{hyperref}
\title{Zul Language Resources for K-12 Curriculum Development}
\author{Compiled by Claude Sonnet 5.5 (claude-sonnet-5-5) for Andr\'as Kornai}
\date{}
\begin{document}
\maketitle

\section{Introduction}
Zul (ISO 639-3: zlu; Glottocode zull1239) is an Afro-Asiatic language of the Chadic branch. It belongs to the Polci cluster of the South Bauchi languages and is spoken in villages in Toro and Bauchi Local Government Areas, Bauchi State, Nigeria. Sources in the search results estimate about 4,000 speakers. Younger speakers are reported to be shifting toward Hausa. Other names include Zule, Zulawa and Mbarmi. A community reading and writing book (2006) and a preliminary Bible translation are mentioned in one summary, but their titles and locations were not found.

\section{Primary Educational Resources}
No primers or textbooks with accessible URLs were found. The 2006 community reading and writing book is mentioned in the Endangered Languages Project summary, but no link to it was seen.

\section{Linguistic and Reference Materials}
\begin{itemize}
\item Glottolog entry: \url{https://glottolog.org/resource/languoid/id/zull1239.bigmap.html}
\item Numerals of Zul, Channumerals database (Max Planck Institute for Evolutionary Anthropology): \url{https://lingweb.eva.mpg.de/channumerals/Zul.htm}
\item South Bauchi languages (Wikipedia): \url{https://en.wikipedia.org/wiki/South_Bauchi_languages}
\item Polci language (Wikipedia), a related language in the same cluster: \url{https://en.wikipedia.org/wiki/Polci_language}
\item An article in the Journal of Linguistics and Language Studies (JOLAN) that came up in the search. Its relevance to Zul was not verified: \url{https://jolan.com.ng/index.php/home/article/download/45/40}
\end{itemize}

\section{Religious and Community Resources}
\begin{itemize}
\item Joshua Project profile of the Zul (Zulle) people. It says the community is split between Christian and Muslim affiliation: \url{https://joshuaproject.net/people_groups/22377}
\end{itemize}

\section{Digital and Audio Resources}
\begin{itemize}
\item Endangered Languages Project entry: \url{https://endangeredlanguages.com/lang/10464?hl=en}
\end{itemize}
No audio or video collections were found.

\section{Curriculum Development Context}
Zul is a small language under pressure from Hausa, and no school use is known. A mother-tongue curriculum would build on the community reading and writing book, if it can be located, and would involve Zul speakers directly in documentation and orthography work. Hausa and English remain the likely languages of wider schooling.

\end{document}